\section{Part C. Modular multiplicities and finite curve quotients}

In \S52.2, folio 744, scans 193--194, Grothendieck asks whether a
connected, smooth, separated algebraic multiplicity of finite type and
relative dimension one over a field $k$ has a presentation
$\mathcal X=[U/G]$, with $U$ a smooth quasi-projective curve and $G$ a
finite group. He immediately reformulates this as the existence of a
finite \'etale scheme cover. Characteristic zero is imposed only after
this question. Folios 745--746, scans 195--197, then discuss arithmetic
and universal moduli multiplicities, including a possible integral
base. We interpret ``multiplicity'' as a Deligne--Mumford stack; a
quotient is a stack quotient, and a finite group means a constant
finite group scheme. These distinctions are essential at every base.

\subsection{The complete characteristic-zero answer}

\begin{statement}{C.1}{Known}{Finite quotient criterion and the classification used}
Let $\mathcal X$ be a connected smooth separated Deligne--Mumford curve
of finite type over $\overline{\mathbb Q}$. The following conditions
are equivalent:
\begin{enumerate}
\item $\mathcal X=[U/G]$ with $U$ a smooth quasi-projective curve and
      $G$ finite;
\item $\mathcal X$ admits a finite \'etale surjective scheme cover;
\item every geometric inertia map
      $I_x\to\pi_1^{\mathrm{et}}(\mathcal X,x)$ is injective;
\item after embedding $\overline{\mathbb Q}$ into $\mathbb C$, the
      analytic universal cover is $\mathbb H$, $\mathbb C$, or
      $\mathbb P^1$, rather than a nontrivial weighted projective stack.
\end{enumerate}
The precise classification is
\citet[Theorem 1.1, Proposition 7.2, Corollary 7.7]{BehrendNoohi2006};
the algebraic inertia criterion and comparison are their Theorems 4.2
and 4.4 and Remark 4.3.
\end{statement}
\begin{proof}
A quotient atlas $U\to[U/G]$ is a finite \'etale $G$-torsor, whether
or not the action is faithful or free on points. Conversely, if
$V\to\mathcal X$ is a degree-$d$ finite \'etale scheme cover, its
frame torsor
$\operatorname{Isom}_{\mathcal X}(\{1,\ldots,d\},V)$ is an
$S_d$-torsor. It is a scheme: it is an open and closed subspace of
the $d$-fold fiber product of $V$ over $\mathcal X$, and the latter
is finite over the scheme $V$. It is a smooth curve, finite over a
quasi-projective curve, hence quasi-projective. This proves (1)--(2)
without an effectiveness hypothesis.

For (2)--(3), stabilizers of a finite \'etale cover are precisely the
inverse images of the corresponding stabilizers in its monodromy set.
There are only finitely many inertia types on a finite-type curve:
one generic type and finitely many jumping points. If all inertia
maps are injective, choose a finite quotient of the profinite group
that separates every nonidentity element of these finitely many
finite groups. Its regular cover has trivial inertia. A smooth
separated algebraic space curve is a scheme, so it gives (2).
This is the curve case of the cited algebraic criterion.

Analytic covering theory gives the analogous criterion for the
universal cover. The classification says that the only simply
connected smooth analytic Deligne--Mumford curves are $\mathbb H$,
$\mathbb C$ and
\[
 \mathbb P(m,n)=[(\mathbb A^2-\{0\})/\mathbb G_m],
 \qquad t\cdot(x,y)=(t^m x,t^n y),\quad m,n\geq1.
\]
Corollary 7.7 supplies a finite curve cover whenever the universal
cover is an ordinary curve. Riemann existence (their Theorem 4.4)
algebraizes this finite analytic cover. The characteristic-zero
invariance of finite \'etale covers under extension of algebraically
closed fields transports it back to $\overline{\mathbb Q}$.
All these are statements about finite covers, not an assertion that
the analytic upper half-plane is an algebraic curve.
\end{proof}

\begin{statement}{C.2}{Refuted}{The exact exceptions for effective orbifolds}
For an effective Deligne--Mumford curve over $\overline{\mathbb Q}$,
the proposed assertion in scan 193 fails precisely for a projective
line with one cone point of order greater than one, or with two cone
points of unequal orders. With the convention that order one means
an ordinary point, these are the footballs $\mathcal F(m,n)$ with
$m\ne n$. Every other effective orbifold curve has a finite curve
quotient presentation. This is
\citet[Proposition 5.7 and Corollary 5.9]{BehrendNoohi2006}.
\end{statement}
\begin{proof}
Put $a=\gcd(m,n)$. Van Kampen gives
\[
 \pi_1^{\mathrm{top}}\mathcal F(m,n)
 =\langle u,v\mid u^m=v^n=uv=1\rangle=C_a.
\]
The inertia maps at the two points are $C_m\to C_a$ and
$C_n\to C_a$, with kernel orders $m/a$ and $n/a$. If $m\ne n$,
at least one kernel is nontrivial. The same calculation is valid
for the \'etale group, since this finite group is already profinite.
C.1 therefore excludes every bad football, including the one-cone
case $n=1$. This is a separate obstruction calculation for each
member of the exceptional family.

When $m=n$, the map $z\mapsto z^m$ gives
$\mathcal F(m,m)=[\mathbb P^1/\mu_m]$. In all remaining signatures,
Proposition 5.7 supplies an ordinary simply connected universal
cover and a finite-index surface subgroup; Corollary 5.9 gives
the finite curve cover. For a noncompact orbifold the familiar
presentation is a free product of a free group and its finite
cyclic inertia groups. Mapping those cyclic factors injectively
into their finite product exhibits a torsion-free finite-index
subgroup directly. The compact higher-genus and triangle cases
are included in the stated theorem, so the exceptional list is
exhaustive.
\end{proof}

\begin{statement}{C.3}{Refuted}{Gerbes contribute further exceptions, with an exact list}
The bad-football list is not the list for arbitrary
Deligne--Mumford curves. Write $\mathcal X\to\mathcal C$ for
rigidification by generic inertia $H$, so that $\mathcal C$ is
effective. All hyperbolic and Euclidean $\mathcal X$ have finite
curve covers. In the spherical case put
$d=|\ker(H\to\pi_1^{\mathrm{top}}\mathcal X)|$.
The exceptional cases are exactly:
\begin{enumerate}
\item $\mathcal C=\mathcal F(m,n)$ and either $m\ne n$ or $d>1$;
\item $\mathcal C$ has cone orders $(2,2,r)$, $r\geq2$,
      $(2,3,3)$, $(2,3,4)$ or $(2,3,5)$, and $d>1$.
\end{enumerate}
The ordinary projective line is included as $\mathcal F(1,1)$.
Their universal covers are, respectively,
$\mathbb P(dm/a,dn/a)$, $a=\gcd(m,n)$, and $\mathbb P(d,d)$.
This is the complete translation of
\citet[Propositions 7.4--7.6 and Corollary 7.7]{BehrendNoohi2006}.
\end{statement}
\begin{proof}
The universal-cover list follows from the cited propositions.
The weighted projective stack $\mathbb P(r,s)$ is an ordinary
curve if and only if $r=s=1$: the stabilizers at its two coordinate
points have orders $r,s$. It is simply connected, because the
homotopy fiber sequence
\[
 \mathbb C^*\longrightarrow\mathbb C^2-\{0\}
          \longrightarrow\mathbb P(r,s)
\]
has $\pi_1(\mathbb C^2-\{0\})=0$ and
$\pi_0(\mathbb C^*)=0$. Thus a nontrivial stabilizer in this
universal cover maps trivially to its fundamental group. Its
image in $\mathcal X$ also lies in the kernel of an inertia map.
C.1 excludes every case listed above. Equivalently, pulling a
putative scheme cover back to this simply connected stack would
give a disjoint union of copies of a nonscheme stack, which cannot
be a scheme. When the listed obstruction is absent, the universal
cover is $\mathbb P^1$ and C.1 applies.

An especially simple missed exception is
$\mathbb P(n,n)$, $n>1$. It is the gerbe of $n$th roots of
$\mathcal O_{\mathbb P^1}(1)$ (changing to $\mathcal O(-1)$
changes only the dual convention). Its generic inertia is
$\mu_n$, whereas its \'etale fundamental group is trivial.
The Kummer class is the residue of $1$ in
$\operatorname{Pic}(\mathbb P^1)/n\cong\mathbb Z/n$,
so the gerbe is nonneutral. It has no finite \'etale scheme cover,
despite having the ordinary projective line as its effective
orbifold. The trivial gerbe $\mathbb P^1\times B\mu_n$, on the
other hand, equals $[\mathbb P^1/\mu_n]$ with trivial action.
This distinguishes a bad gerbe from mere ineffectiveness.

For clarity about the positive theorem, in the nonspherical case
one first takes a finite surface cover of $\mathcal C$ and kills
the finite outer band $\pi_1(C)\to\operatorname{Out}(H)$.
The remaining gerbe class is a torsion class in
$H^2(C,Z(H))$. It vanishes for an open surface. For a compact
surface of positive genus, an unramified cover of degree
$n=\operatorname{exp}Z(H)$ kills it, since pullback on
$H^2$ multiplies by the degree. Such a cover is the kernel of
the surjection $\pi_1(C)\to C_n$ sending one handle generator
to $1$ and the others to $0$. This supplies the finite-cover
step in Corollary 7.7. The source's final displayed list of
generators for that index-$n$ subgroup omits Schreier conjugates;
using the kernel gives the required subgroup without that slip.
Also its Proposition 7.6(ii) prints $k$ in the second weight;
the proof and the cyclic-kernel argument give $d$ in both weights,
as used here.
\end{proof}

\begin{statement}{C.4}{Proved}{Existence of a finite quotient descends to the ground field}
Let $k$ be a characteristic-zero field and $\mathcal X/k$ a
connected smooth separated Deligne--Mumford curve of finite type.
If $\mathcal X_{\bar k}$ admits a finite \'etale scheme cover,
then $\mathcal X=[C/G]$ for a smooth quasi-projective curve $C/k$
and a constant finite group $G$. In particular there is no
additional obstruction to \emph{existence} over $\mathbb Q$.
The curve $C$ can be connected over $k$; it need not be
geometrically connected, and $G$ need not be the group of a
specified presentation over $\bar k$.
\end{statement}
\begin{proof}
A finite \'etale surjective morphism
$U_{\bar k}\to\mathcal X_{\bar k}$, together with the finite
type scheme $U_{\bar k}$, descends to a finite extension $L/k$.
Finiteness, \'etaleness and surjectivity descend after enlarging
$L$: these properties are of finite presentation, and
surjectivity is faithfully flat descent for the resulting
finite map. The composite $U_L\to\mathcal X_L\to\mathcal X$
is finite \'etale and surjective because $L/k$ is finite
separable. As a $k$-scheme, $U_L$ is still a smooth
quasi-projective curve.

Take its frame torsor as in C.1, or a connected Galois closure
in the category of finite \'etale covers of $\mathcal X$.
It is a scheme because it is finite over a fiber product of
copies of $U_L$; finite \'etale descent and normalization give
the same construction. The Galois automorphisms form a
constant finite group $G$ over $k$, and descent identifies
$\mathcal X$ with $[C/G]$. A connected component of a finite
Galois cover, with its stabilizer, is still surjective over
the connected stack, so $C$ may be chosen connected.
This construction deliberately changes the cover and its group.
A possible descent obstruction for one prescribed geometric
presentation does not obstruct the asserted existence.
\end{proof}

\subsection{Integral covers and the elliptic test}

\begin{statement}{C.5}{Proved}{An integral boundary-extension lemma}
Let $\overline{\mathcal Y}$ be a connected smooth separated
Deligne--Mumford stack of finite type over $\mathbb Z$, with a
relative normal-crossings divisor $D$, and let
$\mathcal Y=\overline{\mathcal Y}-D$. Assume every horizontal
irreducible boundary component has a point outside the other
components in each residue characteristic. Then restriction
induces an equivalence of categories of finite \'etale covers
of $\overline{\mathcal Y}$ and $\mathcal Y$.
\end{statement}
\begin{proof}
Normalize $\overline{\mathcal Y}$ in a finite \'etale cover
of $\mathcal Y$; normalization is finite because the base is
excellent. One can make this construction on \'etale scheme
charts and descend it. At the generic point of a boundary
component, take the orders of inertia in a Galois closure,
and call their least common multiple $e$. They are finite
and the characteristic-zero extension is tame.

Fix a prime $p$ and a point of that component alone over
$\mathbb F_p$. The relative Abhyankar lemma,
\citet[Expos\'e XIII, Proposition 5.5]{SGA1}, says that a finite
\'etale cover of the complement whose generic fibers are tame
is, strictly locally there, killed by adjoining roots of the
boundary parameters with exponents prime to $p$. Consequently
$e$ is prime to $p$. Strict localization is unramified at the
generic divisor, so it does not remove its inertia. As the
same argument applies to every prime, $e=1$.

Thus the normalization is unramified at every horizontal
codimension-one point. Every vertical codimension-one point
already belongs to $\mathcal Y$, because $D$ is relative,
so it too is unramified. Purity of the branch locus
\citep[Expos\'e X, Theorem 3.1]{SGA1} makes the normalization
\'etale everywhere. Normality also makes extension of
morphisms unique, proving full faithfulness. The scheme
proof descends along \'etale charts, so nontrivial stack
inertia at small primes does not invalidate the lemma.
\end{proof}

\begin{statement}{C.6}{Proved}{The integral elliptic moduli stack is simply connected}
For the stack of elliptic curves with origin over $\mathbb Z$,
\[
 \pi_1^{\mathrm{et}}(\mathcal M_{1,1,\mathbb Z})=1.
\]
It has no nonempty finite \'etale scheme cover, and in particular
no presentation $[C/G]$ by a scheme and a finite constant group.
This refutes the integral extension of the expectation in
\S52.2, folios 744--746, scans 193--197, even for the principal
one-dimensional moduli example.
\end{statement}
\begin{proof}
We give the integral argument, including residue characteristics
$2$ and $3$. Let $\overline{\mathcal M}_{1,1}$ be the stack of
stable one-pointed genus-one curves. Its boundary consists of
the nodal cubic. Deformation of a node has local equation
$xy=t$, so the boundary is a relative smooth Cartier divisor.
It has a point in every characteristic, for example the
Weierstrass cubic $y^2+xy=x^3$ with origin at infinity:
its tangent cone is $y(y+x)$, with distinct lines in every
characteristic. Thus C.5 identifies the finite \'etale covers
of the open and compactified stacks.

Here is a direct computation for the compactification. Use the
five Weierstrass coefficients $a_1,a_2,a_3,a_4,a_6$, put
\[
\begin{aligned}
b_2&=a_1^2+4a_2,& b_4&=a_1a_3+2a_4,\\
b_6&=a_3^2+4a_6,&
b_8&=a_1^2a_6+4a_2a_6-a_1a_3a_4+a_2a_3^2-a_4^2,\\
c_4&=b_2^2-24b_4,&
\Delta&=-b_2^2b_8-8b_4^3-27b_6^2+9b_2b_4b_6.
\end{aligned}
\]
The open scheme
$W=\mathbb A^5_{\mathbb Z}-V(c_4,\Delta)$ parametrizes
smooth or nodal Weierstrass cubics. Changes of Weierstrass
coordinates form the smooth group scheme
$H=\mathbb G_m\ltimes\mathbb A^3$ with coordinates
$(u,r,s,t)$ and $u$ invertible. The usual substitutions
$x=u^2x'+r$, $y=u^3y'+u^2sx'+t$ give
\[
 \overline{\mathcal M}_{1,1}=[W/H].
\]
To see this presentation integrally, the spaces of sections of
$\mathcal O(2O)$ and $\mathcal O(3O)$ on a smooth or nodal
pointed genus-one curve are locally free of ranks two and
three. Choosing their bases produces $x,y$ and the Weierstrass
equation; changing these bases is exactly the substitution
above. The nodal condition is $\Delta=0$, $c_4\ne0$,
and the smooth condition is $\Delta\ne0$. This argument uses
the general Weierstrass equation, not the short equation
that would invert $2$ and $3$.

The complement of $W$ has codimension at least two.
Indeed, over $\mathbb Q$ the polynomials $c_4,\Delta$ are
coprime (completing the square and cube reduces this to the
coprimality of $A$ and $4A^3+27B^2$). There is no vertical
height-one component, because the coefficient of $a_1^4$ in
$c_4$ is $1$. Purity therefore gives
$\pi_1^{\mathrm{et}}(W)=\pi_1^{\mathrm{et}}(\mathbb A^5_{\mathbb Z})$.
The latter is trivial: a connected finite \'etale cover restricts
over $\mathbb Q$ to a constant number-field extension, since
$\mathbb A^5_{\overline{\mathbb Q}}$ is simply connected by
Riemann existence. Pullback along the zero section extends
that number field to a finite \'etale algebra over $\mathbb Z$.
Minkowski's discriminant bound excludes every nontrivial
number field unramified at all finite primes
\citep[Theorems 4.3 and 4.9]{Milne2020ANT}. Normalization
then identifies the original cover with the trivial one.

Finally let $V\to[W/H]$ be any finite \'etale cover. Its
pullback to $W$ is a disjoint union of copies of $W$. Descent
is a permutation of these copies over
$W\times_{[W/H]}W\cong H\times W$. This scheme is connected,
so the permutation is constant; the identity section and
the unit descent axiom force it to be the identity.
Thus the cover is trivial on $[W/H]$. We have proved
$\pi_1^{\mathrm{et}}(\overline{\mathcal M}_{1,1,\mathbb Z})=1$,
and C.5 proves the displayed formula for the open stack.

The stack has nontrivial geometric inertia already over
$\mathbb Q$: every elliptic curve has $[-1]$. A nonempty
finite \'etale cover of a connected stack with trivial
fundamental group is a disjoint union of copies of the
stack, so it cannot be a scheme in this case. For a proposed
$[C/G]$, the torsor $C\to[C/G]$ would be just such a cover.
This gives the requested contradiction without assuming
that the action on $C$ is faithful.
\end{proof}

\begin{statement}{C.7}{Known}{The same elliptic stack does have finite curve covers over fields}
Over $\mathbb Q$ and $\overline{\mathbb Q}$,
$\mathcal M_{1,1}$ is a finite curve quotient. More explicitly,
the scheme $Y(3)$ of elliptic curves with an ordered basis
of their $3$-torsion gives
\[
 \mathcal M_{1,1,\mathbb Z[1/3]}
       =[Y(3)/\operatorname{GL}_2(\mathbb F_3)].
\]
The group has order $48$; over an algebraically closed field,
fixing the Weil pairing gives a connected component with
group $\operatorname{SL}_2(\mathbb F_3)$ of order $24$.
The group-theoretic level structure and its representability
are constructed explicitly in \citet[\S4; Theorem 7.1]{Bruin2007}.
\end{statement}
\begin{proof}
When $3$ is invertible, $E[3]$ is finite \'etale locally
isomorphic to $(\mathbb Z/3)^2$. Its sheaf of ordered bases
is therefore a finite \'etale
$\operatorname{GL}_2(\mathbb F_3)$-torsor. A basis kills
all automorphisms of an elliptic curve, and the full-level
moduli problem is represented by the smooth affine modular
curve $Y(3)$; the cited explicit construction works over
$\mathbb Z[1/3]$. Forgetting the basis is exactly the stated
quotient. Counting a nonzero first column and a second
column outside its span gives $(9-1)(9-3)=48$; the
determinant map has two values, giving $24$ in its kernel.
The connected components over an algebraically closed field
are distinguished by the two possible primitive Weil
pairings. This cover ramifies or ceases to be \'etale at
the excluded prime $3$, in agreement with C.6; it is not
an integral cover over all of $\mathbb Z$.
\end{proof}

\subsection{The four- and five-pointed genus-zero examples}

\begin{statement}{C.8}{Proved}{Integral fundamental groups for the labelled genus-zero moduli}
For the fine moduli schemes of labelled pointed curves,
\[
 \pi_1^{\mathrm{et}}(M_{0,4,\mathbb Z})=1,
 \qquad
 \pi_1^{\mathrm{et}}(M_{0,5,\mathbb Z})=1.
\]
Here $M_{0,4}$ has relative dimension one, and $M_{0,5}$
has relative dimension two.
\end{statement}
\begin{proof}
Normalize the first three marks to $\infty,0,1$. Then
\[
 M_{0,4}=\operatorname{Spec}\mathbb Z
       [t,t^{-1},(t-1)^{-1}],
\]
and $M_{0,5}$ is the open subscheme of $\mathbb A^2_{\mathbb Z}$
with $t,u,t-1,u-1,t-u$ invertible. The compactifications are
$\overline M_{0,4}=\mathbb P^1_{\mathbb Z}$ and
\[
 \overline M_{0,5}
  =\operatorname{Bl}_{(0,0),(1,1),(\infty,\infty)}
       (\mathbb P^1\times\mathbb P^1)_{\mathbb Z}.
\]
The second formula follows by resolving the three triple
collisions of the two moving marks with the fixed marks;
the exceptional coordinate records the third component's
cross-ratio. It also constructs the universal stable
five-pointed curve, so this is the stable compactification.

The boundaries are relative normal-crossings divisors.
For $n=4$ they are the three sections. For $n=5$ they are
the strict transforms of the six horizontal/vertical lines
at $0,1,\infty$, the diagonal, and the three exceptional
divisors. Every boundary component's open stratum has a
geometric point in every characteristic. C.5 therefore
reduces the asserted fundamental groups to those of the
compactifications.

Over $\overline{\mathbb Q}$ both compactifications are
simply connected: $\mathbb P^1\times\mathbb P^1$ is so,
and blowing up a smooth point of a complex surface does
not change its fundamental group. One can see the latter
by van Kampen: the boundary of a small ball is $S^3$, and
the replacing disc bundle over $\mathbb P^1$ is simply
connected. Riemann existence gives the \'etale assertion.
Both compactifications have $\mathbb Z$-sections (boundary
sections suffice). For any finite \'etale cover, its
generic fiber is thus a constant number-field extension,
and pullback to this section makes that extension unramified
at every finite prime. Minkowski gives the trivial field,
and normality gives the trivial cover, exactly as in C.6.
\end{proof}

\begin{statement}{C.9}{Proved}{Symmetric-group quotients retain their stack inertia}
Let $\mathcal U_4=[M_{0,4}/S_4]$ and
$\mathcal U_5=[M_{0,5}/S_5]$ over $\mathbb Z$. Then
\[
 \pi_1^{\mathrm{et}}(\mathcal U_4)=S_4,
 \qquad
 \pi_1^{\mathrm{et}}(\mathcal U_5)=S_5.
\]
Both are global quotients by the displayed smooth
quasi-projective scheme. For $n=4$ the action has kernel
$V_4$, and the effective quotient has
$\pi_1^{\mathrm{et}}([M_{0,4}/S_3])=S_3$ instead.
For $n=5$ the permutation action is faithful.
$\mathcal U_5$ is a quotient of a surface, not a curve;
its dimension prevents a finite curve presentation.
\end{statement}
\begin{proof}
For a connected scheme $X$ and constant finite group $G$,
the quotient atlas is a finite \'etale $G$-torsor and gives
the covering exact sequence
\[
 1\longrightarrow\pi_1^{\mathrm{et}}(X)
 \longrightarrow\pi_1^{\mathrm{et}}([X/G])
 \longrightarrow G\longrightarrow1.
\]
It does not require a free action on geometric points:
freeness is on the torsor over the \emph{stack}. Apply C.8.
This remains true when a residue characteristic divides
$|G|$, because a constant finite group scheme is \'etale
in every characteristic. Tameness of inertia and \'etaleness
of this atlas are different conditions.

Permuting four marks acts on the cross-ratio through the
six rational functions
\[
 t,\quad 1-t,\quad 1/t,\quad 1/(1-t),\quad
 t/(t-1),\quad (t-1)/t.
\]
The kernel consists of the identity and the three double
transpositions, namely $V_4$. The program below computes
all $24$ permutations, with multiplicity four for each
function. Retaining $S_4$ therefore retains a generic
$V_4$-gerbe; replacing it by $S_3$ changes the stack and
its fundamental group.

For five marks faithfulness can already be checked on
boundary divisors of the stable compactification. They
are indexed by two-element subsets of a five-element
set; a permutation acting trivially on the open scheme
acts trivially on its compactification and hence on all
these subsets. A permutation fixing every two-element
subset fixes every individual element, so is the identity.
Finally a finite surjective morphism preserves relative
dimension; since $\dim M_{0,5}/\mathbb Z=2$, no curve can
be its finite atlas. The two-dimensional example is
included because the question specifically asks for it,
not as a one-dimensional instance of C.1.
\end{proof}

\begin{statement}{C.10}{Proved}{Finite checks and their limits}
The program \texttt{code/c\_stack\_checks.py} checks all
$144$ ordered football pairs $1\leq m,n\leq12$, all $24$
cross-ratio permutations, the $48$ invertible $2\times2$
matrices over $\mathbb F_3$, the Weierstrass polynomial
coprimality calculation, and the integral nodal example.
These finite checks verify the numerical and algebraic
ingredients used above. The classification, descent,
purity and fundamental-group conclusions are proved in
C.1--C.9 and are not inferred from sampling.
\end{statement}
\begin{proof}
The source and output record the exact bounds and each
test. Football fundamental-group and inertia-kernel
orders are computed by greatest common divisors, as in
C.2. The cross-ratio computation substitutes all
permutations of $(\infty,0,1,t)$ and simplifies over
$\mathbb Q(t)$; the matrix enumeration tests all $3^4$
matrices. Symbolic calculation verifies
$\gcd_{\mathbb Q[a_1,a_2,a_3,a_4,a_6]}(c_4,\Delta)=1$
and $c_4=1$, $\Delta=0$ on $y^2+xy=x^3$.
The derivative and tangent-cone calculation at its node
is integral, so the conclusion in characteristics $2$
and $3$ does not depend on testing finitely many primes.
\end{proof}
